\section{Insert a Node in a BST}
\label{sec:bst-insert}

\problemstatement{Given the root of a BST and a value $X$ to insert,
insert $X$ into the tree so that the BST invariant still holds, and
return the (possibly new) root.}

\begin{patternbox}
\emph{``Insert into a BST''} reuses Section~\ref{sec:bst-search}'s
search walk almost unchanged: insertion is what happens when a search for
$X$ runs off the edge of the tree (reaches a null pointer) instead of
finding a match -- that null position is exactly where $X$ belongs.
\end{patternbox}

\subsection*{Understanding the problem carefully}

Walk down the tree exactly as in Section~\ref{sec:bst-search}'s search:
at each node, go left if $X$ is smaller, right if $X$ is larger. Since
$X$ is not yet in the tree, this walk will never find an exact match --
it will instead reach a null child pointer. That null position is the
unique place where $X$ can be attached while preserving the BST
invariant: it is smaller than everything on the path so far that went
right, and larger than everything that went left, by construction of the
walk itself.

\begin{ideabox}[Same Walk as Search, Attach at the Null]
Walk from the root as in a search for $X$. Whenever the walk would move
into a null child, create a new node with value $X$ and attach it there
instead. If the tree was empty to begin with, the new node simply becomes
the root.
\end{ideabox}

\subsection*{Why attaching at the first null reached is always valid}

Every step of the walk maintains an invariant: $X$ is consistent with
becoming a descendant of the current node, given everything decided so
far (it is greater than every ancestor we went right from, and less than
every ancestor we went left from). Reaching a null pointer means no
existing node currently occupies the position consistent with all those
decisions -- so creating a new node there extends the tree without
violating the BST invariant anywhere, since the new node's relationship
to every ancestor on the path is already exactly correct by construction.

\subsection*{Algorithm}

\begin{enumerate}
  \item If the current node is null: create and return a new node with
        value $X$.
  \item If $X <$ current node's value: recursively insert into the left
        child, and attach the (possibly updated) result back as the left
        child.
  \item Else: recursively insert into the right child, attach back as
        the right child.
  \item Return the current node (unchanged, except possibly for the
        newly attached child).
\end{enumerate}

\begin{algorithm}[H]
\caption{Insert into a BST}
\begin{algorithmic}[1]
\Procedure{InsertBST}{\textit{root}, $X$}
  \If{$\textit{root} = \text{null}$}
    \State \Return new node with value $X$
  \EndIf
  \If{$X < \textit{root}.\textit{val}$}
    \State $\textit{root}.\textit{left} \gets \Call{InsertBST}{\textit{root}.\textit{left}, X}$
  \Else
    \State $\textit{root}.\textit{right} \gets \Call{InsertBST}{\textit{root}.\textit{right}, X}$
  \EndIf
  \State \Return \textit{root}
\EndProcedure
\end{algorithmic}
\end{algorithm}

\subsection*{Dry run}

\dryrun{Insert $X=6$ into the BST with root $5$, left child $3$, right
child $8$.}

\begin{itemize}
  \item At $5$: $6>5$, go right to $8$.
  \item At $8$: $6<8$, go left -- but $8$ has no left child (null).
        Create new node $6$, attach as $8$'s left child.
\end{itemize}

The tree now has $8$'s left child equal to $6$.

\subsection*{C++ implementation}

\begin{lstlisting}
#include <bits/stdc++.h>
using namespace std;

struct TreeNode {
    int val;
    TreeNode* left;
    TreeNode* right;
    TreeNode(int v) : val(v), left(nullptr), right(nullptr) {}
};

TreeNode* insertBST(TreeNode* root, int X) {
    if (root == nullptr) return new TreeNode(X);

    if (X < root->val) root->left = insertBST(root->left, X);
    else root->right = insertBST(root->right, X);

    return root;
}

int main() {
    TreeNode* root = new TreeNode(5);
    root->left = new TreeNode(3);
    root->right = new TreeNode(8);

    insertBST(root, 6);
    cout << "8's left child: " << root->right->left->val << endl;
    return 0;
}
\end{lstlisting}

\begin{complexitybox}
\textbf{Time Complexity.} $O(h)$, identical to search.
\textbf{Space Complexity.} $O(h)$ for the recursion stack, plus $O(1)$
for the single new node created.
\end{complexitybox}

\begin{takeawaybox}
Insertion is search that creates a node instead of reporting failure --
recognizing this means you do not need to separately re-derive the walk
logic, only add the ``attach a new node here'' step at the point where
search would otherwise give up.
\end{takeawaybox}
